从实验结果可以看出，Padding Oracle 攻击并不是去破解 AES 算法本身，而是利用了解密系统在实现上的信息泄露。攻击者虽然不知道密钥，但只要能够反复提交构造密文，并区分服务器返回的是“填充正确”还是“填充错误”，就可以逐步恢复 CBC 解密过程中的中间状态，最后还原出真实明文。

\subsection{实验结果分析}

本实验能够成功，依赖于以下几个条件同时成立：
\begin{itemize}
    \item 目标系统使用了 CBC 分组模式，并采用 PKCS\#7 填充；
    \item 解密服务器在解密后会检查填充是否合法；
    \item 攻击者能够区分服务器返回的不同错误状态。
\end{itemize}

这三个条件一旦同时满足，服务器实际上就成了一个“padding oracle”。虽然它没有直接把明文或密钥返回出来，但每一次反馈都在泄露一部分关于解密结果的信息。攻击者把这些信息不断累积起来，就能一点一点地恢复出完整明文。

\subsection{实现过程中的体会}

这次实验中比较明显的一点是：原理本身并不算太难，真正容易出错的是实现细节。例如，填充长度的判断、字节下标的对应关系、中间状态与真实明文的区分，这些地方只要错一个，后面的恢复过程就会全部受到影响。

\subsection{安全启示}

本实验说明：只有机密性、没有完整性保护的 CBC 加密实现是不安全的。为了防御 Padding Oracle 攻击，实际系统中通常要避免把可区分的 padding 错误直接暴露给外部用户，更安全的做法包括：
\begin{itemize}
    \item 使用带认证的加密模式，如 GCM、CCM 等；
    \item 在解密前先做消息认证，而不是先解密再判断 padding；
    \item 对外统一错误返回，避免泄露“是 padding 错误还是其他错误”这类差异信息。
\end{itemize}

总的来说，这次实验不仅加深了对 CBC 模式和 PKCS\#7 填充规则的理解，也更直观地说明了：密码算法本身安全，并不等于系统实现就一定安全。
